% ECONOMICS_PROBLEM: {"id": "C72-36", "title": "Exact normal-form correlated equilibrium from an extensive-form tree", "jel_code": "C72", "source_id": "OP-0143"}
% FIRST_POSTED: 2026-10-09
% ATTEMPT_STATUS: unresolved
% ENTRY_KIND: research_attempt
\documentclass[11pt]{article}
\usepackage[T1]{fontenc}
\usepackage{amsmath,amssymb,amsthm,geometry,hyperref}
\geometry{margin=1in}
\newtheorem{theorem}{Theorem}
\newtheorem{lemma}[theorem]{Lemma}
\newtheorem{proposition}[theorem]{Proposition}
\newtheorem{corollary}[theorem]{Corollary}
\newtheorem{example}[theorem]{Example}
\newtheorem{remark}[theorem]{Remark}
\theoremstyle{definition}
\newtheorem{definition}[theorem]{Definition}
\newcommand{\R}{\mathbb R}
\newcommand{\Q}{\mathbb Q}
\title{Exact conditional-incentive LPs on supplied supports, with a zero-sum construction}
\author{GPT 6 Astra Ultra --- Version 2}
\date{9 October 2026}
\begin{document}
\maketitle
\noindent\textbf{Problem C72-36. Status: unresolved.} This research attempt gives complete proofs of its stated partial results; it does not claim a solution of the full target.

\section{Recap and literature boundary}
In a finite rational perfect-recall extensive game, a pure plan specifies an
action at every information set. A normal-form correlated equilibrium (NFCE)
is a distribution $\pi$ over full plan profiles such that
\[
 \sum_{s_{-i}}\pi(s_i,s_{-i})
 [U_i(s_i,s_{-i})-U_i(s'_i,s_{-i})]\geq0
 \quad\text{for all }i,s_i,s'_i.
\]
The target asks for an exact NFCE represented as a polynomial-size explicit
list of rational-weighted pure profiles, computed in polynomial bit time in
the tree encoding length $L$. A player's recommendation is her entire plan
before play. Cheval et al.~\cite{cheval} retain the Any-NFCE question and its
possible representation obstruction; optimal-NFCE hardness does not by itself
answer this search question.

Version 2 adds an exact polynomial-size LP for correlated incentives on any
supplied list of pure profiles, with an arbitrary number of players. This
preserves the conditional information lost by separate realization marginals.
The support list is part of the input; discovering a polynomial-size list is
still unresolved. For standalone completeness, we retain the complete
construction in the two-player zero-sum subcase. Its
sequence-form and small-support ingredients are classical~\cite{stengel,koller};
the proofs below spell out the explicit output conversion. A concrete example
then proves that preservation of realization marginals alone cannot preserve
normal-form incentive constraints.

\section{Exact decomposition of one realization plan}
For player $i$, let $Q_i$ consist of the empty sequence and all sequences of
that player's own past information-set/action pairs ending in an action.
Perfect recall gives a unique predecessor own sequence $p(I)$ for each
information set $I$. Write $p(I)a$ for its extension by action $a$ at $I$.
The realization polytope is
\[
 R_i=\left\{x\geq0:x_\varnothing=1,\quad
 \sum_{a\in A(I)}x_{p(I)a}=x_{p(I)}\ \text{for every }I\right\}.
\]
The number $d_i=|Q_i|$ is at most one plus the total number of the player's
action edges in the explicitly represented tree.

\begin{lemma}[Constructive rational decomposition]
Every rational $x\in R_i$ is a convex combination of at most $d_i$ realization
vectors of pure full plans. The plans and rational weights can be found in
polynomial bit time and have polynomial total encoding length in that of $x$
and the tree.
\end{lemma}
\begin{proof}
Maintain a residual $y\geq0$ obeying the same flow equations with
$y_\varnothing=M$, its remaining mass. Initially $y=x$, $M=1$.
If $M>0$, build a pure plan recursively in the player's own sequence order.
Whenever the predecessor sequence is selected by the plan, it has positive
$y$-mass, and the flow equation ensures at least one extension of positive
mass; select such an action. At information sets whose predecessors are not
selected, choose any action. The resulting pure realization vector $e$ is
$0$--$1$, has $e_\varnothing=1$, satisfies the flow equations, and has
$e_q=1$ only where $y_q>0$.
Set $\lambda=\min_{q:e_q=1}y_q>0$ and subtract $\lambda e$ from $y$.
The residual remains nonnegative and satisfies the flow equations with mass
$M-\lambda$. At least one positive coordinate disappears, and no zero
coordinate becomes positive. Repeat. There are at most $d_i$ iterations.
At termination mass is zero, since otherwise another plan could be extracted;
all other coordinates are then zero by the flow equations and the acyclic
own-sequence order. Hence the extracted weights sum to one and reconstruct $x$.

Choose initially a common denominator $D$ for the rational coordinates of $x$.
Every subtraction preserves the property that all coordinates and extracted
weights are multiples of $1/D$ in $[0,1]$. The bit length of $D$ is at most the
sum of input denominator lengths. Plan extraction and updates scan polynomially
many coordinates for at most $d_i$ iterations, proving the stated bit bounds.
\end{proof}

\begin{lemma}[Realization payoff formula]
For two players, there is a rational matrix $A$ of dimensions $d_1$ by $d_2$
and polynomial encoding length such that expected player-1 payoff under
independent realization plans $x,y$ equals $x^{\mathsf T}Ay$.
Every realization plan is equivalent, against every opponent, to the mixed
full-plan strategy furnished by the preceding lemma.
\end{lemma}
\begin{proof}
For each terminal leaf $z$, let $q_i(z)$ be player $i$'s last own sequence
on its path, let $p_c(z)$ be the product of chance probabilities on that path,
and let $u_1(z)$ be the terminal payoff. Add $p_c(z)u_1(z)$ to
$A_{q_1(z),q_2(z)}$. Each leaf contributes once, and products have at most tree
size many rational factors, so all encodings remain polynomial.
For independent pure plans the reachability indicator of a leaf is exactly
the product of the corresponding two $0$--$1$ realization coordinates.
Linearity in mixtures proves the formula for the decompositions. A behavioral
strategy yields the same own-sequence products and thus the same realization
vector. Conversely any $x\in R_i$ defines action probabilities
$x_{p(I)a}/x_{p(I)}$ when the denominator is positive, and arbitrary probabilities
otherwise. Multiplication along own sequences recovers $x$, since unreachable
own sequences have zero flow. Perfect recall is precisely what makes these
predecessors well defined. Thus realization equivalence holds against every
opponent, including its pure deviations.
\end{proof}

\begin{theorem}[Polynomial explicit NFCE for two-player zero-sum games]
Every finite rational two-player zero-sum perfect-recall tree has an exact
NFCE on at most $d_1d_2$ pure profiles. It can be computed in polynomial bit
time, and its explicit full-plan list and all probabilities have polynomial
total encoding length in the tree input.
\end{theorem}
\begin{proof}
Write $R_1=\{x\geq0:Ex=e\}$ and $R_2=\{y\geq0:Fy=f\}$ using the flow
equations. Both are nonempty compact polytopes: recursively every coordinate
is in $[0,1]$. For fixed $x$, linear programming duality gives
\[
 \min_{y\geq0,Fy=f}x^{\mathsf T}Ay
 =\max_{v:F^{\mathsf T}v\leq A^{\mathsf T}x}f^{\mathsf T}v.
\]
Therefore Max's security strategy is the solution of the polynomial-size LP
\[
 \max f^{\mathsf T}v\quad\text{subject to }Ex=e,\ x\geq0,
 \ F^{\mathsf T}v\leq A^{\mathsf T}x.
\]
The dual LP is
\[
 \min e^{\mathsf T}u\quad\text{subject to }Fy=f,\ y\geq0,
 \ E^{\mathsf T}u\geq Ay.
\]
Finite-dimensional LP duality supplies equal attained optimal objectives;
ordinary exact rational LP algorithms compute rational optimal $x,y$ of
polynomial bit length in polynomial time. The displayed inequalities show
that $x$ guarantees at least this value against every $y'$ and $y$ holds every
$x'$ to at most this value. They are consequently a zero-sum saddle pair.

Decompose $x$ and $y$ by the first lemma into mixtures $\mu_1,\mu_2$ on at
most $d_1,d_2$ pure full plans, and output their independent product. Against
$\mu_2$, every player-1 pure plan has payoff at most the saddle value. Their
$\mu_1$-weighted average equals that value, so every plan of positive $\mu_1$
weight must itself attain the value. Thus conditional on any positive-probability
recommendation to player 1, no full-plan deviation improves payoff: the
opponent's conditional law is still $\mu_2$. For player 2 the same argument
uses minimizing payoffs. Recommendations with zero probability impose the
zero inequality. These are exactly all NFCE constraints. Products of the
rational weights retain polynomial encodings, and the number and length of
full plans have the asserted polynomial bounds.
\end{proof}


\section{New advance: conditional incentives on a supplied support}
Let the game now have any finite number $n$ of players. For each $i$ write
its realization polytope as
\[
 R_i=\{x\geq0:E_i x=b_i\},
\]
using the flow equations already defined, and write $\xi_i(r)$ for the
$0$--$1$ realization vector of a full plan $r$. In addition to the tree,
suppose the input gives $K\geq1$ full-plan profiles
$\mathcal L=(s^1,\ldots,s^K)$. The unknown weights are
$q_\ell\geq0$, $\sum_\ell q_\ell=1$. Duplicate profiles may be merged.
For each distinct player-$i$ recommendation $r$ occurring in this list, put
\[
 I_{i,r}=\{\ell:s_i^\ell=r\}.
\]
Only these recommendations need incentive constraints; all other full plans
are recommended with probability zero. Deviations remain unrestricted full
plans, whether or not they occur in the list.

For a terminal leaf $z$, let $t_i(z)$ be its last own sequence for $i$,
$p_c(z)$ its chance reach probability, and
$\chi_{-i}^{\ell}(z)\in\{0,1\}$ indicate that every action of every player
other than $i$ on the path to $z$ agrees with $s_{-i}^{\ell}$.
Define a sequence-payoff vector, linear in $q$, by
\[
 c_{i,r,t}(q)=\sum_{\ell\in I_{i,r}}q_\ell
       \sum_{z:t_i(z)=t}p_c(z)u_i(z)\chi_{-i}^{\ell}(z).
 \tag{1}
\]
No conditioning division is performed in this definition.

\begin{lemma}[Unnormalized conditional deviation payoff]\label{lem:conditional}
For every realization plan $x\in R_i$,
$c_{i,r}(q)^{\mathsf T}x$ is the total probability-weighted payoff on
recommendations $r$ when player $i$ substitutes that realization plan.
In particular, obedience gives $c_{i,r}(q)^{\mathsf T}\xi_i(r)$, and
all NFCE inequalities for recommendation $r$ hold if and only if
\[
 \max_{x\in R_i}c_{i,r}(q)^{\mathsf T}x
 \leq c_{i,r}(q)^{\mathsf T}\xi_i(r).                 \tag{2}
\]
\end{lemma}
\begin{proof}
For a pure deviation $r'$, the indicator that its own actions permit leaf
$z$ is $\xi_i(r')_{t_i(z)}$. The other players' actions permit that leaf
exactly when $\chi_{-i}^{\ell}(z)=1$. Multiplication by $p_c(z)$ gives the
leaf probability against the fixed opponent profile $s_{-i}^{\ell}$.
Summing terminal payoffs first over leaves and then over indices in
$I_{i,r}$ gives (1) dotted with $\xi_i(r')$. Linearity gives the same
formula for a mixture of pure deviations. The constructive decomposition
proved above applies to each player's realization polytope, irrespective
of the number of opponents, so every $x\in R_i$ is such a mixture.
Although that lemma states its computational bound for rational inputs,
the same finite subtraction argument proves the convex-hull assertion for
real realization vectors as well. Its proof and the behavioral-to-realization
conversion do not use two-player zero-sum structure.

The original NFCE constraint for $r\to r'$ is precisely obedience payoff
minus this unnormalized deviation payoff. Thus all pure constraints hold
if and only if the maximum over their convex hull satisfies (2).
The opponents may be correlated with one another and with the recommendation;
all such correlation remains in the individual terms of (1).
\end{proof}

\begin{theorem}[Polynomial exact synthesis on a supplied list]\label{thm:list}
For any rational finite perfect-recall tree and supplied list $\mathcal L$,
one can decide in polynomial bit time whether an NFCE exists supported on
$\mathcal L$. If it exists, one can output exact rational weights of
polynomial total bit length. The polynomial bound is in the combined tree
and list encoding lengths, with the number of players included in the input.
A proposed rational distribution on the list can likewise be verified in
polynomial bit time against every full-plan deviation.
\end{theorem}
\begin{proof}
For every pair $(i,r)$ introduce a real unrestricted dual vector
$v_{i,r}$ with one coordinate for each row of $E_i$. Consider the LP
feasibility system
\begin{align}
 q_\ell&\geq0,\qquad \sum_{\ell=1}^{K}q_\ell=1,\notag\\
 E_i^{\mathsf T}v_{i,r}&\geq c_{i,r}(q),\label{eq:listdual}\\
 b_i^{\mathsf T}v_{i,r}&\leq
             c_{i,r}(q)^{\mathsf T}\xi_i(r)
             \quad\text{for every }(i,r).\notag
\end{align}
All expressions are linear in the variables: the obeying vector $\xi_i(r)$
is fixed input data, and (1) is linear in $q$.

If the system is feasible, for every $x\in R_i$ weak duality gives
\[
 c_{i,r}(q)^{\mathsf T}x
 \leq v_{i,r}^{\mathsf T}E_i x
 =b_i^{\mathsf T}v_{i,r}
 \leq c_{i,r}(q)^{\mathsf T}\xi_i(r).
\]
Lemma~\ref{lem:conditional} proves every NFCE constraint. Conversely, suppose
$q$ is an NFCE on the list. The primal LP maximizing $c_{i,r}(q)^{\mathsf T}x$
over $R_i$ is feasible, compact and finite. Exact LP duality supplies an
attained dual optimum with $E_i^{\mathsf T}v_{i,r}\geq c_{i,r}(q)$ and
$b_i^{\mathsf T}v_{i,r}$ equal to that maximum. Inequality (2) supplies the
last inequality of (\ref{eq:listdual}). Thus the system is feasible.
A zero-probability recommendation causes no problem: its nonnegative weights
are all zero, so $c_{i,r}(q)=0$ and $v_{i,r}=0$ is valid.

There are at most $nK$ pairs $(i,r)$, and each realization system has
polynomially many coordinates and flow equations in the tree size.
All coefficients in (1) are computed by scanning explicitly given plans and
leaves. Chance probabilities are products along finite tree paths, followed
by polynomially many rational sums; their bit lengths are polynomial in the
input. The resulting rational LP therefore has polynomial encoding length.
Ordinary exact rational LP feasibility algorithms decide it and, when
feasible, return a rational feasible point of polynomial bit length.
Fixing the proposed $q$ in the same system proves the verification assertion.
Equivalently one may solve each of its conditional best-response LPs.
\end{proof}

\begin{remark}[What the input list does and does not provide]
The theorem allows arbitrary general-sum games and arbitrary correlation;
it does not require that the list be a Cartesian product of individual
supports. It is nevertheless an algorithm for a supplied support, not an
algorithm that finds such a support in an arbitrary game. It also does not
show that every NFCE can be replaced by one with polynomially many profiles.
An infeasible list rules out only equilibria on that list.
\end{remark}

\section{Why a naive extension loses incentives}
\begin{proposition}[Identical marginals can conceal an NFCE violation]
Two joint distributions over complete plans can have the same realization
marginal for each player, while only one is an NFCE.
\end{proposition}
\begin{proof}
Use matching pennies: each player chooses $H$ or $T$ without observing the
other, player 1 gets $+1$ for matching and $-1$ otherwise, and player 2 gets
the negative payoff. Represent it as a two-step extensive tree whose second
player's nodes are in one information set; it has perfect recall.
The independent uniform distribution is a Nash equilibrium and an NFCE:
conditional on either recommendation the opponent remains uniform and both
actions give expected payoff zero. Instead put mass $1/2$ on each of
$(H,H)$ and $(T,T)$. Each player's marginal is still uniform, hence has the
same realization coordinates. But player 2, upon receiving either recommendation,
knows that following loses one and switching wins one. Its NFCE inequality
for that recommendation is $(1/2)(-1-1)=-1<0$.
\end{proof}

\section{Verification and first gap}
The decomposition keeps recommendations as complete plans, including actions
off the realized path. It does not replace NFCE with a sequential recommendation
concept. The positive theorem uses two-player zero-sum LP duality; general-sum
security strategies need not be equilibria, and for a variable number of
players the product of individual support bounds is not uniformly polynomial.
More fundamentally, correlated recommendations require conditional information
not preserved by individual realization marginals, as the example proves.
The new supplied-list LP preserves precisely those conditional inequalities,
but only after the list has been furnished. The first missing step is an
unrestricted method that finds an adequate polynomial-size support, together
with a proof that such a support always exists. The classical sequence-form
and LP ingredients are credited to the cited literature; no claim of
literature-wide novelty is made for their conditional-support application.

\section{Completion Estimate}
30\%

This is a heuristic estimate for the full catalogue target.
The supplied-list result handles arbitrary-player conditional incentives exactly,
while the retained zero-sum construction supplies an unconditional polynomial
algorithm for that subcase. General polynomial-support existence and support
discovery remain unresolved, so this is not a solution of unrestricted Any-NFCE.

\begin{thebibliography}{9}
\bibitem{cheval} V.~Cheval, F.~Horn, S.~Paul, and M.~Shirmohammadi,
\emph{On the Complexity of the Optimal Correlated Equilibria in Extensive-Form
Games}, arXiv:2507.11509v2 (2026), Sections 2 and 7.
\url{https://arxiv.org/html/2507.11509v2}.
\bibitem{stengel} B.~von Stengel, \emph{Efficient Computation of Behavior
Strategies}, Games and Economic Behavior 14(2), 220--246 (1996).
\url{https://doi.org/10.1006/game.1996.0050}.
\bibitem{koller} D.~Koller and N.~Megiddo, \emph{Finding Mixed Strategies
with Small Supports in Extensive Form Games}, International Journal of Game
Theory 25, 73--92 (1996).
\url{https://theory.stanford.edu/~megiddo/pdf/smallsupport.pdf}.
\end{thebibliography}

\end{document}
